\documentclass[12pt]{article}
\usepackage{geometry}
\usepackage{titlesec}
\usepackage{setspace}
\usepackage{hyperref}
\usepackage{enumitem}
\usepackage{lmodern}
\usepackage{microtype}

\geometry{margin=1in}
\setstretch{1.25}

% Section title formatting
\titleformat{\section}
  {\large\bfseries\filcenter}
  {}{0em}{}
\titlespacing*{\section}{0pt}{1.2em}{0.6em}

\titleformat{\subsection}
  {\normalsize\bfseries}
  {}{0em}{}
\titlespacing*{\subsection}{0pt}{1em}{0.4em}

%-----------------------------
\begin{document}

\begin{center}
    {\LARGE \bfseries Exploring the Mathematical and Physical Foundations of Relativity} \\[0.4em]
    {\normalsize Anant Krishna and Sujoge Dua — QLab (Period 5)} \\[0.25em]
    {\small Research Proposal — October 2025}
\end{center}

\vspace{1em}

%----------------------------------------
\section*{Introduction}

This project aims to explore and understand the mathematical and physical foundations of Einstein’s theories of special and general relativity. Rather than beginning with a fixed research question, the goal is to develop a strong conceptual and mathematical grounding in relativity and use that foundation to identify a specific topic for deeper investigation later in the year. 

Our study will emphasize the connections between geometry, algebra, and physics in describing spacetime. We will begin with special relativity, covering Lorentz transformations, spacetime intervals, and the Minkowski metric, and gradually extend to general relativity, focusing on curvature and tensor analysis. By doing so, we aim to gain a comprehensive understanding of how mathematical structures encode the behavior of space, time, and gravity.

The broader value of this project lies in its emphasis on deep theoretical understanding before specialization. This approach mirrors the way physicists build research questions from first principles, allowing us to approach advanced topics in relativity with both rigor and clarity.

%----------------------------------------
\section*{Specific Goals and Measurements}

\textbf{Goal 1 — Foundational Understanding.}  
Develop a detailed understanding of the mathematical and conceptual framework underlying both special and general relativity. This includes deriving and interpreting Lorentz transformations, understanding spacetime geometry, and studying how tensors describe curvature and gravitation. Progress will be measured by the clarity and completeness of derivations, summaries, and notes demonstrating comprehension of these core topics.

\textbf{Goal 2 — Topic Identification and Research Design.}  
After building the necessary background, identify a focused topic within relativity—such as gravitational time dilation, geodesic motion, or relativistic electromagnetism—for deeper exploration. Success will be measured by the formulation of a specific research question and a clear plan for investigating it in the second phase of the project.

%----------------------------------------
\section*{Specific Responsibilities}

For the most part, we both share the same responsibilities as we are both required to learn the topic.

%----------------------------------------
\section*{Potential Pitfalls and Alternative Strategies}

\textbf{Challenge 1 — Mathematical Complexity.}  
The tensor calculus and differential geometry used in general relativity are advanced topics.  
\textit{Strategy:} begin with the mathematical framework of special relativity, emphasizing conceptual understanding and geometric intuition before tackling the full tensor formulation.

\textbf{Challenge 2 — Scope Refinement.}  
Relativity is a vast field, and finding a manageable research focus may be challenging.  
\textit{Strategy:} maintain flexibility in early phases, using readings and derivations to naturally identify an area that is both interesting and feasible to study in depth. We can also consult Dr. Linch for guidance.

%----------------------------------------
\section*{Proposed Timeline}

\textbf{October – December 2025: Foundational Phase}  
Study the core principles of special relativity and the mathematical tools needed for general relativity through Dr. Linch's packets and guidance.
\textit{Goal:} develop a rigorous understanding of spacetime structure and prepare for topic selection.

\textbf{January – March 2026: Exploration Phase}  
Begin studying general relativity and survey possible research directions (e.g., gravitational effects, relativistic motion, or spacetime curvature).
\textit{Goal:} identify and formally define a specific topic for focused study.

\textbf{March – May 2026: Research Phase}  
Conduct deeper analysis or theoretical work on the chosen topic, culminating in a final report or presentation.  
\textit{Goal:} synthesize understanding and produce clear written and visual explanations.

\textbf{Paper Writing Timeline}  
\begin{itemize}[leftmargin=2em, itemsep=0.2em]
    \item December 2025 — draft foundational background and theory.
    \item March 2026 — define specific topic and begin structured research.
    \item May 2026 — complete final paper and presentation.
\end{itemize}

%----------------------------------------
\section*{Annotated Citations}

Einstein, A. (1905). “Zur Elektrodynamik bewegter Körper.” \textit{Annalen der Physik, 17}, 891–921.  
Introduces special relativity and reformulates time and space in terms of invariant light speed. Establishes the foundation for Lorentz transformations and the geometry of spacetime.

Minkowski, H. (1908). “Raum und Zeit.” \textit{Physikalische Zeitschrift, 10}, 75–88.  
Proposes the unification of space and time into a four-dimensional spacetime manifold, forming the mathematical basis for modern relativity.

Einstein, A. (1915). “Die Feldgleichungen der Gravitation.” \textit{Sitzungsberichte der Königlich Preußischen Akademie der Wissenschaften}, 844–847.  
Presents the Einstein field equations, linking the curvature of spacetime to the energy–momentum content of matter, which will inform the general relativity portion of this study.

\end{document}
